% 第 11 题程序框图（含「输入」框——本项目首例）：按固定源 PDF 矢量路径逐框重绘（1 pt → 0.03515 cm）。
% 原点取主干竖线（源页 x 564.62）与「结束」底边（y 688.94），y 向上。共 8 框 8 箭头，全实线。
% 结构要点（易错）：「输入 a,b」与「输出 i」同为向右倾平行四边形；i=i+1 与 a=a+i,b=b-i 同高并列于右侧；
% 「否」支自菱形右顶点水平右行、折向上入 i=i+1 底边；i=i+1 的出口在其【上边】，向上至 y=118.60 后
% 折左汇入主干——汇入点在 i=1 与 a=a+i,b=b-i 之间，故循环体含处理框整段（非尾循环）。
% 判断框内为单线大于号（ASCII >），非斜线双划。
\begin{tikzpicture}[>=stealth,every node/.style={font=\small},x=0.03515cm,y=0.03515cm]
% ===== 框 =====
\draw[rounded corners=3.98pt] (-13.99,186.03) rectangle (13.99,202.87);   % 开始
\draw (-30.20,154.37)--(23.97,154.37)--(30.19,171.46)--(-23.98,171.46)--cycle; % 输入（右倾平行四边形）
\draw (-14.93,125.89) rectangle (14.93,139.80);   % i=1
\draw (-47.20,95.01) rectangle (47.20,111.32);    % a=a+i,b=b-i
\draw (61.76,95.76) rectangle (108.06,110.56);    % i=i+1（右侧同高）
\draw (0,80.36)--(34.81,71.72)--(0,63.08)--(-34.81,71.72)--cycle; % 菱形 a>b
\draw (-23.69,31.57)--(17.54,31.57)--(23.68,48.43)--(-17.55,48.43)--cycle; % 输出（右倾平行四边形）
\draw[rounded corners=3.98pt] (-13.99,0) rectangle (13.99,17.01);          % 结束
% ===== 框内文字 =====
\node at (0,194.45) {开始};
\node at (0,162.92) {输入 \(a,b\)};
\node at (0,132.85) {\(i=1\)};
\node at (0,103.17) {\(a=a+i,b=b-i\)};
\node at (84.91,103.16) {\(i=i+1\)};
\node at (0,71.72) {\(a>b\)};
\node at (0,40.00) {输出 \(i\)};
\node at (0,8.51) {结束};
% ===== 箭头（主干六段 + 否支 + 回接）=====
\draw[->] (0,185.83)--(0,171.46);   % 开始→输入
\draw[->] (0,154.37)--(0,139.80);   % 输入→i=1
\draw[->] (0,125.89)--(0,111.32);   % i=1→处理框
\draw[->] (0,95.01)--(0,80.36);     % 处理框→菱形
\draw[->] (0,63.08)--(0,48.43);     % 菱形→输出（是分支）
\draw[->] (0,31.57)--(0,17.01);     % 输出→结束
\draw[->] (34.81,71.72)--(84.91,71.72)--(84.91,95.76);  % 否支：右顶点→右→上入 i=i+1 底边
\draw[->] (84.91,110.56)--(84.91,118.60)--(0,118.60);   % 回接：i=i+1 上边→上→左汇入主干
% ===== 分支标注 =====
\node at (9.13,56.29) {是};
\node at (35.32,63.87) {否};
\end{tikzpicture}
